\documentclass[aspectratio=169]{beamer}

% -------------------- Packages --------------------
\usepackage{graphicx}
\usepackage{amsmath}
\usepackage{booktabs}
\usepackage{adjustbox}
\usepackage{amssymb}
\usepackage{pifont}

% -------------------- Theme --------------------
\usetheme{Madrid}
\usecolortheme{default}
\setbeamertemplate{caption}[numbered]

% -------- Attractive Accent Color (Teal) --------
\definecolor{structurecolor}{RGB}{22,160,133}
\setbeamercolor{structure}{fg=structurecolor}
\setbeamercolor{frametitle}{bg=structurecolor, fg=white}

% -------------------- CUSTOM FOOTER --------------------
\setbeamertemplate{footline}{}
\setbeamertemplate{navigation symbols}{}
\setbeamertemplate{footline}
{
  \leavevmode
  \hbox{
  \begin{beamercolorbox}[
    wd=\paperwidth,
    ht=3ex,
    dp=1.2ex,
    leftskip=0.4cm,
    rightskip=0.4cm
  ]{custom footer}
  \color{white}\scriptsize
  Lekshmi R Rajan\;|\; Roll No.: 31
  \hfill  \insertframenumber/\inserttotalframenumber
  \end{beamercolorbox}}
}
\setbeamercolor{custom footer}{bg=structurecolor}

% -------------------- Title Info --------------------
\title{AI-Enabled Home Planning and 3D Visualization Web Application}
\author[Lekshmi R Rajan | Roll No.: 31]{
\textbf{Lekshmi R Rajan}\\
Roll No.: 31\\
Reg. No.: KTE24MCA--2035
}
\institute{
\textbf{Guided By:}\\
Prof. Shilpa M Thomas\\[0.2cm]
Department of Computer Applications\\
Rajiv Gandhi Institute of Technology (RIT), Kottayam
}
\date{}

\begin{document}

%---------------- TITLE ----------------
\begin{frame}
\titlepage
\end{frame}

%---------------- INTRODUCTION ----------------
\begin{frame}{INTRODUCTION}
\begin{itemize}
  \item Home planning requires a clear understanding of spatial layout, room functionality, and user requirements.
  \item During the early planning stage, hand-sketched house plans are commonly used to express layout ideas.
  \item These hand-drawn sketches are difficult for non-technical users to interpret and visualize accurately.
  \item Existing architectural design tools (e.g., CAD software) demand professional expertise and are not user-friendly for common users.
  \item The proposed AI-enabled web-based system converts hand-sketched house plans into interactive 3D visualizations via a full-stack pipeline using PyTorch, OpenCV, FastAPI, Three.js, and Groq LLM.
\end{itemize}
\end{frame}

%---------------- OBJECTIVE ----------------
\begin{frame}{OBJECTIVE}
\begin{itemize}
  \item To design a secure, authenticated web interface for uploading hand-sketched house plans.
  \item To apply a hybrid AI pipeline — PyTorch U-Net for wall segmentation and OpenCV for geometry extraction — for sketch interpretation.
  \item To generate an interactive 3D visualization of residential house layouts using Three.js (WebGL).
  \item To clearly visualize functional room areas (Living Room, Bedroom, Kitchen, Bathroom) with color differentiation.
  \item To estimate approximate construction cost using an ensemble of 4 Scikit-Learn ML models (Linear Regression, Decision Tree, Random Forest, Gradient Boosting).
  \item To provide context-aware, page-specific design and budget advice via an AI Chatbot powered by Groq LLM (\texttt{llama-3.1-8b-instant}).
\end{itemize}
\end{frame}

%---------------- SCOPE ----------------
\begin{frame}{SCOPE}
\textbf{Project Includes:}
\begin{itemize}
  \item Secure user authentication (JWT Bearer Tokens + Argon2 password hashing).
  \item AI-powered upload and processing of hand-sketched house plans.
  \item Generation of a conceptual 3D house layout with 5-room standard structure.
  \item Visualization of functional areas (Living, Bedroom, Kitchen, Bathroom) via Three.js rendering.
  \item Approximate construction cost estimation (Economy, Standard, Premium tiers) using ML.
  \item Context-aware Groq LLM chatbot with strict Design vs.\ Budget mode separation.
\end{itemize}
\vspace{0.3cm}
\textbf{Project Excludes:}
\begin{itemize}
  \item Detailed structural engineering calculations.
  \item Construction-ready blueprints (mechanical/electrical/plumbing).
  \item Replacement of certified architectural sign-offs.
\end{itemize}
\end{frame}

%---------------- RELEVANCE ----------------
\begin{frame}{RELEVANCE}
\begin{itemize}
  \item \textbf{Democratizes Design:} Enables easy visual understanding for non-technical users without CAD expertise.
  \item \textbf{Cost \& Time Efficient:} Reduces sketch-to-3D prototyping from hours to $\sim$2.6 seconds (99\% time reduction).
  \item \textbf{Decision Support:} Real-time ML cost estimates enable confident financial planning at the ideation stage.
  \item \textbf{Technology Integration:} Demonstrates a unified pipeline of Deep Learning (U-Net), Classical CV (OpenCV), ML Regression (Scikit-Learn), and LLM (Groq) in a single web system.
  \item \textbf{Practical Utility:} Supports early-stage budgeting and spatial planning for homeowners and architects alike.
\end{itemize}
\end{frame}

%---------------- REQUIREMENT ANALYSIS ----------------
\begin{frame}
\centering
\Huge \textbf{REQUIREMENT ANALYSIS}
\end{frame}

%---------------- EXISTING SYSTEM ----------------
\begin{frame}{EXISTING SYSTEM (LITERATURE REVIEW)}
\scriptsize
\begin{table}
\centering
\begin{tabular}{p{3.2cm} p{2.4cm} p{3.2cm} p{4cm}}
\toprule
\textbf{Publication / Paper Title} & \textbf{Authors \& Year} & \textbf{Methodology} & \textbf{Key Limitation} \\
\midrule
Automatic Floor Plan Generation &
Liu et al. (2022) &
Computer Vision + CNN &
Converts sketches to vector plans; no 3D output \\
3D House Modeling using CAD &
Autodesk Research (2021) &
Manual CAD Modeling &
Accurate but requires professional expertise \\
Web-based 3D Interior Visualization &
Kim et al. (2023) &
Three.js + WebGL &
Interactive 3D but requires manual data entry \\
AI-based Cost Estimation &
Patel et al. (2022) &
ML Regression &
Predicts cost but lacks visual integration \\
\bottomrule
\end{tabular}
\caption{Literature Review: Existing Home Planning Systems}
\end{table}
\end{frame}

%---------------- GAP IDENTIFIED ----------------
\begin{frame}{GAP IDENTIFIED}
\begin{itemize}
  \item \textbf{No Integrated Solution:} No single platform combines Sketch Upload $\rightarrow$ AI Segmentation $\rightarrow$ 3D Model $\rightarrow$ Cost Estimation.
  \item \textbf{Steep Learning Curve:} Existing CAD tools are too complex for average users with no technical background.
  \item \textbf{Static 2D Limitations:} Hand sketches fail to convey spatial volume, depth, and room flow.
  \item \textbf{Budget Opacity:} Lack of data-driven, real-time initial cost estimation integrated with design tools.
  \item \textbf{No Conversational Assistance:} No system provides spatially-grounded LLM chatbot linked to the actual extracted layout data.
\end{itemize}
\end{frame}

%---------------- PROPOSED SYSTEM ----------------
\begin{frame}{PROPOSED SYSTEM}
\begin{itemize}
  \item The proposed system is a full-stack AI-enabled web application that accepts a hand-sketched house plan as input and automatically converts it into an interactive 3D model.
  \item A hybrid AI pipeline first preprocesses the sketch using OpenCV (CLAHE, Canny Edge Detection, Hough Transform) and then applies a trained PyTorch U-Net model for pixel-level wall segmentation.
  \item The extracted wall geometry is mapped to real-world metre coordinates and rendered interactively in the browser using Three.js (WebGL) with rotate, zoom, and pan controls.
  \item A Scikit-Learn ensemble of 4 ML regression models (Linear Regression, Decision Tree, Random Forest, Gradient Boosting) automatically estimates construction cost from the extracted layout dimensions.
  \item A Groq LLM chatbot (\texttt{llama-3.1-8b-instant}) provides context-aware advice — acting as a Design Assistant on the Visualization page and as a Cost Consultant on the Cost Estimation page.
  \item User accounts, project data, and uploaded sketches are securely managed via JWT authentication, Argon2 password hashing, and SQLite persistence.
\end{itemize}
\end{frame}

%---------------- SOFTWARE & HARDWARE ----------------
\begin{frame}{SOFTWARE \& HARDWARE REQUIREMENT}
\begin{columns}[T]
\begin{column}{0.48\textwidth}
\textbf{Software Requirements}
\begin{itemize}
  \item \textbf{Backend:} Python 3.10, FastAPI, Uvicorn
  \item \textbf{Frontend:} HTML5, CSS3, JavaScript (Three.js, OrbitControls.js)
  \item \textbf{Deep Learning:} PyTorch (U-Net model)
  \item \textbf{CV:} OpenCV (cv2)
  \item \textbf{ML:} Scikit-Learn (4 regressors)
  \item \textbf{LLM:} Groq SDK (\texttt{llama-3.1-8b-instant})
  \item \textbf{Auth:} Passlib (Argon2), JWT (python-jose)
  \item \textbf{DB:} SQLite via SQLModel
\end{itemize}
\end{column}
\begin{column}{0.48\textwidth}
\textbf{Hardware Requirements}
\begin{itemize}
  \item \textbf{Processor:} Intel Core i5 or higher
  \item \textbf{RAM:} Minimum 8 GB (16 GB recommended)
  \item \textbf{GPU:} NVIDIA GPU recommended for U-Net inference; browser WebGL for 3D rendering
  \item \textbf{Storage:} 500 MB minimum (96 MB model + uploads)
  \item \textbf{Internet:} Required for Groq API calls
\end{itemize}
\end{column}
\end{columns}
\end{frame}

%---------------- DEVELOPMENT METHODOLOGY ----------------
\begin{frame}{DEVELOPMENT METHODOLOGY}
\begin{itemize}
\item \textbf{Agile Approach:} Iterative development with continuous testing and integration.
\item \textbf{Sprint 1 (Research \& Planning):} Problem definition, literature survey, requirement analysis, technology selection.
\item \textbf{Sprint 2 (Backend \& AI Pipeline):} FastAPI setup, OpenCV preprocessing module, U-Net training and inference, cost model training (4 ML algorithms).
\item \textbf{Sprint 3 (Frontend \& 3D):} Three.js 3D viewer with OrbitControls, room coloring, interactive UI with upload, visualization, and cost pages.
\item \textbf{Sprint 4 (Integration):} JWT authentication, project persistence (SQLite/SQLModel), Groq LLM chatbot integration with dual-mode context prompting.
\item \textbf{Sprint 5 (Testing \& Validation):} Unit testing, end-to-end pipeline validation, latency benchmarking, and final documentation.
\end{itemize}
\end{frame}

%---------------- SYSTEM ARCHITECTURE DIAGRAM ----------------
\begin{frame}{SYSTEM ARCHITECTURE OVERVIEW}
\begin{figure}
\centering
\includegraphics[height=0.75\textheight]{l1.jpeg}
\caption{High-Level System Architecture of the AI Home Planner}
\end{figure}
\end{frame}

%---------------- 3-TIER ARCHITECTURE ----------------
\begin{frame}{SYSTEM ARCHITECTURE: 3-TIER INFRASTRUCTURE}
\begin{itemize}
    \item \textbf{Tier 1 — Presentation Layer:} HTML5/CSS3/JavaScript frontend with Three.js for 3D rendering. Pages: Home, Upload, Visualization, Cost Estimation, Dashboard.
    \item \textbf{Tier 2 — Application Layer:} FastAPI + Uvicorn backend serving REST APIs. Manages JWT authentication, file uploads (UUID-named), and asynchronous subprocess execution via Python's \texttt{subprocess} module.
    \item \textbf{Tier 3 — AI \& Data Layer:} PyTorch U-Net for wall segmentation $\rightarrow$ OpenCV geometry extraction $\rightarrow$ Scikit-Learn cost prediction. Groq LLM handles conversational queries. SQLite (via SQLModel) stores users and project data.
    \item \textbf{Data Flow:} Upload $\rightarrow$ \texttt{layout\_extraction.py} subprocess $\rightarrow$ \texttt{layout\_to\_3d.py} subprocess $\rightarrow$ JSON response to Three.js renderer.
\end{itemize}
\end{frame}

%---------------- ARCHITECTURE FIGURES ----------------
\begin{frame}{SYSTEM ARCHITECTURE: PREPROCESSING PIPELINE}
\begin{figure}
\centering
\includegraphics[height=0.7\textheight]{l2.jpeg}
\caption{Computer Vision Preprocessing \& Geometry Extraction (OpenCV)}
\end{figure}
\end{frame}

\begin{frame}{SYSTEM ARCHITECTURE: DEEP LEARNING MODULE}
\begin{figure}
\centering
\includegraphics[height=0.7\textheight]{l3.jpeg}
\caption{Deep Learning Semantic Segmentation — PyTorch U-Net Architecture}
\end{figure}
\end{frame}

\begin{frame}{SYSTEM ARCHITECTURE: ML COST ENGINE}
\begin{figure}
\centering
\includegraphics[height=0.7\textheight]{l4.jpeg}
\caption{Machine Learning Cost Estimation Ensemble (Scikit-Learn)}
\end{figure}
\end{frame}

\begin{frame}{SYSTEM ARCHITECTURE: GROQ LLM CHATBOT}
\begin{figure}
\centering
\includegraphics[height=0.7\textheight]{l5.jpeg}
\caption{Groq LLM Context Engine — Dual-Mode Chatbot Architecture}
\end{figure}
\end{frame}

\begin{frame}{SYSTEM ARCHITECTURE: FASTAPI APPLICATION TIER}
\begin{figure}
\centering
\includegraphics[height=0.7\textheight]{l6.jpeg}
\caption{FastAPI Application Tier: Asynchronous Orchestration \& JWT Security}
\end{figure}
\end{frame}

%---------------- IMPLEMENTATION ALGORITHMS ----------------
\begin{frame}{IMPLEMENTATION ALGORITHMS (1/3): Computer Vision}
\textbf{OpenCV Preprocessing Pipeline}
\begin{itemize}
    \item \textbf{Grayscale Conversion:} Reduces input to single-channel image for processing efficiency.
    \item \textbf{CLAHE (Contrast Limited Adaptive Histogram Equalization):} Enhances local contrast on faint, lightly-penciled sketches.
    \item \textbf{Canny Edge Detection:} Multi-stage gradient algorithm that identifies sharp structural boundary lines.
    \item \textbf{Probabilistic Hough Transform:} Connects fragmented or broken wall line segments into continuous lines.
    \item \textbf{Suzuki's Contour Detection Algorithm:} Extracts closed room boundaries from the binary wall mask.
\end{itemize}
\end{frame}

\begin{frame}{IMPLEMENTATION ALGORITHMS (2/3): AI \& ML}
\textbf{Deep Learning \& Machine Learning Modules}
\begin{itemize}
    \item \textbf{PyTorch U-Net:} Pixel-level semantic segmentation of walls vs.\ background. Model: \texttt{wall\_segmentation\_model.pth} (96 MB). Uses skip connections to preserve fine wall boundary details.
    \item \textbf{Groq LLM (\texttt{llama-3.1-8b-instant}):} Context-aware chatbot with \emph{strict} page-based system prompting. Visualization page $\rightarrow$ Design Mode. Cost page $\rightarrow$ Budget Mode. Returns structured JSON responses with optional \texttt{action} objects (e.g., \texttt{change\_color}, \texttt{apply\_theme}, \texttt{add\_furniture}).
    \item \textbf{4 Cost Models (Scikit-Learn):} Linear Regression (baseline), Decision Tree, Random Forest ($R^2$=0.9983, best), Gradient Boosting.
\end{itemize}
\end{frame}

\begin{frame}{IMPLEMENTATION ALGORITHMS (3/3): Security \& 3D}
\textbf{Security \& 3D Rendering}
\begin{itemize}
    \item \textbf{Argon2 Password Hashing:} Industry-standard memory-hard hash (via Passlib) for secure user password storage in SQLite.
    \item \textbf{JWT Bearer Token Authentication:} All protected endpoints (\texttt{/projects/}, \texttt{/users/me/}) require valid OAuth2 access tokens generated on login.
    \item \textbf{UUID File Naming:} Every uploaded file is assigned a unique \texttt{uuid4} prefix to prevent overwrite conflicts.
    \item \textbf{Three.js Linear Extrusion:} 2D room bounding boxes (in metres) are extruded along the Z-axis to a height of 3.0m to create 3D wall meshes.
    \item \textbf{Raycasting (Three.js):} Detects user mouse-click interactions on 3D room surfaces for room selection and labeling.
\end{itemize}
\end{frame}

%---------------- TOOLS & TECHNOLOGIES ----------------
\begin{frame}{TOOLS \& TECHNOLOGIES USED}
\textbf{Development Environment:} VS Code, Git / GitHub, Anaconda / pip

\vspace{0.2cm}
\textbf{Frontend Stack:}
\begin{itemize}
    \item HTML5, CSS3, JavaScript (ES6+)
    \item Three.js \& OrbitControls.js (3D rendering engine)
\end{itemize}

\vspace{0.2cm}
\textbf{Backend \& AI Stack:}
\begin{itemize}
    \item FastAPI + Uvicorn (REST API server)
    \item OpenCV (cv2) — Image preprocessing
    \item PyTorch — U-Net deep learning inference
    \item Scikit-Learn — ML cost regression (joblib model persistence)
    \item Groq SDK — \texttt{llama-3.1-8b-instant} LLM chatbot
    \item SQLite + SQLModel — User \& project database persistence
    \item Passlib (Argon2) + python-jose (JWT) — Authentication
\end{itemize}
\end{frame}

%---------------- CURRENT STATUS: 100% ----------------
\begin{frame}{PROJECT STATUS: 100\% IMPLEMENTATION COMPLETE}
\textbf{All Modules Fully Implemented \& Tested:}
\begin{itemize}
    \item[\ding{51}] \textbf{Secure Auth:} User registration, JWT login, Argon2 password hashing -- all endpoints protected.
    \item[\ding{51}] \textbf{Upload \& Processing:} Image upload with UUID naming $\rightarrow$ U-Net DL segmentation $\rightarrow$ OpenCV fallback $\rightarrow$ layout JSON via subprocess.
    \item[\ding{51}] \textbf{3D Visualization:} Three.js viewer renders 5-room layout with real metre dimensions, room coloring, rotate/zoom/pan controls.
    \item[\ding{51}] \textbf{Cost Estimation:} 4 trained ML models predict cost; Random Forest selected as best ($R^2 = 0.9983$).
    \item[\ding{51}] \textbf{Groq LLM Chatbot:} Dual-mode chatbot live -- Design Mode (Visualization page) and Budget Mode (Cost page).
    \item[\ding{51}] \textbf{Project Persistence:} All projects saved to SQLite and retrievable from the Dashboard.
\end{itemize}
\end{frame}

%---------------- RESULTS: INTERFACE ----------------
\begin{frame}{RESULTS: PLATFORM INTERFACE}
\begin{columns}[T]
\begin{column}{0.33\textwidth}
\begin{figure}
\centering
\includegraphics[width=0.9\linewidth]{fig1.png}
\caption{Home Page \& Landing Interface}
\end{figure}
\end{column}
\begin{column}{0.33\textwidth}
\begin{figure}
\centering
\includegraphics[width=0.9\linewidth]{fig4.png}
\caption{Upload \& AI Processing UI}
\end{figure}
\end{column}
\begin{column}{0.33\textwidth}
\begin{figure}
\centering
\includegraphics[width=0.9\linewidth]{fig5.png}
\caption{Interactive 3D Visualization}
\end{figure}
\end{column}
\end{columns}
\end{frame}

%---------------- RESULTS: DASHBOARD ----------------
\begin{frame}{RESULTS: DASHBOARD \& COST OUTPUT}
\begin{columns}[T]
\begin{column}{0.33\textwidth}
\begin{figure}
\centering
\includegraphics[width=0.9\linewidth]{fig6.png}
\caption{3D Layout: Top-Down View}
\end{figure}
\end{column}
\begin{column}{0.33\textwidth}
\begin{figure}
\centering
\includegraphics[width=0.9\linewidth]{fig7.png}
\caption{ML Cost Estimation Report}
\end{figure}
\end{column}
\begin{column}{0.33\textwidth}
\begin{figure}
\centering
\includegraphics[width=0.9\linewidth]{fig8.png}
\caption{Responsive Dashboard View}
\end{figure}
\end{column}
\end{columns}
\end{frame}

%---------------- RESULTS: MODEL EVALUATION ----------------
\begin{frame}{RESULTS: ML MODEL EVALUATION}
\begin{columns}[T]
\begin{column}{0.48\textwidth}
\begin{figure}
\centering
\includegraphics[width=0.9\linewidth]{fig9.png}
\caption{U-Net Evaluation: mIoU 0.884, Pixel Acc 95.2\%}
\end{figure}
\end{column}
\begin{column}{0.48\textwidth}
\begin{figure}
\centering
\includegraphics[width=0.9\linewidth]{fig10.png}
\caption{Cost Model Comparison: Random Forest Best (MAE \$1,575)}
\end{figure}
\end{column}
\end{columns}
\end{frame}

%---------------- PROJECT PLAN ----------------
\begin{frame}{PROJECT PLAN \& TIMELINE}
\scriptsize
\begin{table}
\centering
\begin{tabular}{p{2cm} p{6.5cm} p{3.5cm}}
\toprule
\textbf{Phase} & \textbf{Activities} & \textbf{Status} \\
\midrule
Phase 1 & Literature Survey, Problem Definition \& Technology Selection & \textbf{Completed} (Dec 2025) \\
Phase 2 & Requirement Analysis \& System Design & \textbf{Completed} (Jan 2026) \\
Phase 3 & Backend: FastAPI, U-Net, OpenCV, ML Cost Model Training & \textbf{Completed} (Jan 2026) \\
Phase 4 & Frontend: Three.js 3D Viewer, UI Pages, JWT Auth & \textbf{Completed} (Feb 2026) \\
Phase 5 & AI Integration: Groq LLM Chatbot, Full Pipeline Testing & \textbf{Completed} (Feb 2026) \\
Phase 6 & Validation, Benchmarking \& Final Documentation & \textbf{Completed} (Mar 2026) \\
\bottomrule
\end{tabular}
\caption{Project Phases and Timeline — 100\% Complete}
\end{table}
\end{frame}

%---------------- CONCLUSION ----------------
\begin{frame}{CONCLUSION}
\begin{itemize}
  \item The AI Home Planner was successfully designed and implemented as a complete, full-stack web application that bridges the gap between hand-drawn sketches and professional 3D architectural visualization.
  \item The hybrid AI pipeline combining OpenCV preprocessing and PyTorch U-Net semantic segmentation effectively interprets noisy, informal sketches with high accuracy (mIoU: 0.884, Pixel Accuracy: 95.2\%).
  \item The Scikit-Learn cost estimation engine (Random Forest: $R^2 = 0.9983$, MAE = \$1,575.64) enables users to receive instant, reliable construction budget estimates derived automatically from the extracted layout.
  \item The Groq LLM chatbot (\texttt{llama-3.1-8b-instant}) provides spatially-grounded, page-specific responses -- acting as a Design Assistant on the Visualization page and a Cost Consultant on the Cost Estimation page.
  \item The complete ``Upload-to-3D'' pipeline runs in approximately 2.6 seconds, representing a 99\% reduction in time compared to traditional manual 3D modeling.
  \item The system successfully democratizes architectural prototyping, making early-stage home planning accessible, visual, and data-driven for both non-technical homeowners and professional architects.
\end{itemize}
\end{frame}

%---------------- GIT HISTORY ----------------
\begin{frame}{GIT VERSION CONTROL HISTORY}
\begin{figure}
\centering
\includegraphics[height=0.7\textheight]{git.png}
\caption{Git Commit History — Full Development Lifecycle}
\end{figure}
\end{frame}

%---------------- BIBLIOGRAPHY ----------------
\begin{frame}{BIBLIOGRAPHY}
\small
\begin{thebibliography}{9}
\bibitem{ref1}
O. Ronneberger, P. Fischer, and T. Brox,
``U-Net: Convolutional Networks for Biomedical Image Segmentation,''
\textit{MICCAI}, pp. 234--241, 2015.
\bibitem{ref2}
G. Bradski, ``The OpenCV Library,''
\textit{Dr. Dobb's Journal}, 2000.
\bibitem{ref3}
F. Pedregosa et al., ``Scikit-learn: Machine Learning in Python,''
\textit{JMLR}, vol. 12, pp. 2825--2830, 2011.
\bibitem{ref4}
R. Cabello, ``Three.js -- JavaScript 3D Library,''
[Online]. Available: https://threejs.org/. 2026.
\bibitem{ref5}
Patel et al., ``ML Approaches for Construction Cost Estimation,''
\textit{Journal of Construction Engineering}, 2022.
\end{thebibliography}
\end{frame}

%---------------- THANK YOU ----------------
\begin{frame}
\centering
\Huge \textbf{THANK YOU}
\vspace{0.5cm} \\
\Large \textbf{Questions \& Answers}
\end{frame}

\end{document}
